\section{Ish va energiya}

\begin{enumerate}

\item Minorali kran uzunligi 5 m va kesimi 0,01 $m^2$ bo‘lgan po‘lat balkani gorizontal holatida 15 m balandlikka tekis ko‘taradi. Kran bajaradigan ishni toping. Balka materialining zichligi 7800 $kg/m^3$. $g = 10 m/s^2$.
\begin{enumerate}[label=\Alph*)]
\item 29250
\item 5850
\item 58500
\item 117000
\end{enumerate}

\item 2 kg massali jismni 3 $m/s^2$ tezlanish bilan 1 m balandlikka ko‘tarishda kishi qanday ish bajaradi? $g = 10 m/s^2$.
\begin{enumerate}[label=\Alph*)]
\item 26
\item 20
\item 13
\item 32
\end{enumerate}

\item Yukni doimiy tezlanish bilan vertikal yuqoriga ko‘tara boshlashadi. Harakatning birinchi sekundida bajarilgan ish keyingi, ikkinchi sekundida bajarilgan ishdan necha marta kichik?
\begin{enumerate}[label=\Alph*)]
\item 2
\item 3
\item 4
\item 5
\end{enumerate}

\item 1 kg massali yukni ip orqali vertikal yuqoriga o‘zgarmas tezlanish bilan ko‘tarishadi. 2 s ichida arqonning taranglik kuchi bilan 48 J ish bajarildi. Yukning tezlanishini toping. $g = 10 m/s^2$.
\begin{enumerate}[label=\Alph*)]
\item 4
\item 1
\item 5
\item 2
\end{enumerate}

\item 20 N kuch 2 kg massali yukni 5 $m/s^2$ tezlanish bilan qiya tekislik bo‘ylab 2,5 m balandlikka ko‘tarib qanday ish bajaradi? Kuch qiya tekislikka parallel ta’sir qiladi. Ishqalanish inobatga olinmasin. $g = 10 m/s^2$.
\begin{enumerate}[label=\Alph*)]
\item 100
\item 50
\item 200
\item 150
\end{enumerate}

\item 20 kg massali jismni qiya tekislik bo‘ylab 6 m balandlikka ko‘tarishadi, bunda u tekislik bo‘ylab 10 m o‘tadi. Agar tortish kuchi tekislikka parallel va ishqalanish koeffitsiyenti 0,2 bo‘lsa, ishqalanish kuchining ishini (modulini) toping. $g = 10 m/s^2$.
\begin{enumerate}[label=\Alph*)]
\item 160
\item 640
\item 320
\item 400
\end{enumerate}

\item Massasi 0,5 kg bo‘lgan jism 7 m balandlikdagi qiya tekislik uchidan uning asosigacha sirpanib tushadi. Tekislikning qiyalik burchagi gorizontga $45^\circ$, va ishqalanish koeffitsiyenti 0,2. Ishqalanish kuchining ishini (modulini) toping. $g = 10 m/s^2$.
\begin{enumerate}[label=\Alph*)]
\item 14
\item 3,5
\item 7
\item 10
\end{enumerate}




\section{Kinitik energiya va ish}

\item Jism qandaydir balandlikdan 10 m/s tezlik bilan gorizontal otilgan. Necha sekunddan keyin jismning kinetik energiyasi ikki marta ortadi? $g = 10 m/s^2$.
\begin{enumerate}[label=\Alph*)]
\item 2
\item 1
\item 0,5
\item 1,5
\end{enumerate}

\item Qandaydir balandlikdan 20 m/s tezlik bilan tosh gorizontal otildi. Otishdan 4 s o‘tib toshning kinetik energiyasi 3000 J ga teng bo‘ldi. Toshning massasi qanday? $g = 10 m/s^2$.
\begin{enumerate}[label=\Alph*)]
\item 3
\item 1,5
\item 6
\item 4,5
\end{enumerate}

\item 3 kg massali jism yer sirtidan gorizontga $60^\circ$ burchak ostida 8 m/s tezlik bilan otildi. Jismning ko‘tarilishning eng yuqori nuqtasidagi kinetik energiyasini toping.
\begin{enumerate}[label=\Alph*)]
\item 12
\item 48
\item 36
\item 24
\end{enumerate}

\item Toshning maksimal ko‘tarilish nuqtasidagi kinetik energiyasi otish nuqtasidagi kinetik energiyasining 25\% ni tashkil qilishi uchun toshni gorizontga qanday burchak ostida otish kerak? Javobni graduslarda bering.
\begin{enumerate}[label=\Alph*)]
\item 30
\item 45
\item 60
\item 75
\end{enumerate}

\item Avtomobil tekis tezlanuvchan harakat boshladi. Harakatning dastlabki 10 s da uning kinetik energiyasining o‘zgarishi keyingi 10 s dagi o‘zgarishidan necha marta kichik?
\begin{enumerate}[label=\Alph*)]
\item 3
\item 2
\item 4
\item 9
\end{enumerate}

\item 0,8 kg massali basketbol to‘pi 10 m/s tezlik bilan uchadi. O‘yinchi to‘pni tutib oladi va 0,1 s ichida uni to‘xtatadi. O‘yinchi qanday o‘rtacha quvvatga erishadi? (Quvvatning absolyut qiymatini ko‘rsating).
\begin{enumerate}[label=\Alph*)]
\item 200
\item 400
\item 800
\item 600
\end{enumerate}

\item 4 kg massali erkin tushayotgan jismning tezligi qandaydir masofada 2 m/s dan 8 m/s gacha oshdi. Shu yo‘ldagi og‘irlik kuchining ishini toping.
\begin{enumerate}[label=\Alph*)]
\item 60
\item 240
\item 120
\item 180
\end{enumerate}

\item 200 g massali tosh gorizontal yuzadan gorizontga burchak ostida otildi va shu yuzaga 40 m masofada 4 s dan so‘ng kelib tushdi. Shu otishga sarflangan ish qanchaga teng? $g = 10 m/s^2$.
\begin{enumerate}[label=\Alph*)]
\item 25
\item 50
\item 100
\item 75
\end{enumerate}

\item 800 m/s tezlik bilan gorizontal uchayotgan 5 g massali o‘q taxtani teshib o‘tadi va undan 400 m/s tezlikda uchib chiqadi. Taxtaning qarshilik kuchi tomonidan o‘q ustida bajarilgan ishning absolyut qiymatini toping.
\begin{enumerate}[label=\Alph*)]
\item 1200
\item 600
\item 2400
\item 1800
\end{enumerate}

\item Qandaydir tezlik bilan uchayotgan o‘q qum joylangan qopga tegadi va uning ichida 15 cm chuqurlikkacha kirib ketadi. Agar shuncha massali o‘qning harakat tezligi ikki marta katta bo‘lsa, o‘q qumga qanday chuqurlikkacha (cm da) kirib ketadi? Qum ichida o‘qning harakatiga qarshilik kuchi uning tezligiga bog‘liq emas deb hisoblang.
\begin{enumerate}[label=\Alph*)]
\item 30
\item 45
\item 75
\item 60
\end{enumerate}





\section{Potensial energiya va ish}

\item 30 m balandlikdagi minoradan tosh gorizontal otildi. Harakat boshlanganidan 2 s o‘tgandagi toshning potensial energiyasini toping. Toshning massasi 0,2 kg. Yer yuzasidagi potensial energiya nolga teng. $g = 10 m/s^2$.
\begin{enumerate}[label=\Alph*)]
\item 40
\item 20
\item 60
\item 10
\end{enumerate}

\item 2 kg massali jism yer sirtidan 6 m/s tezlik bilan gorizontga $30^\circ$ burchak ostida otildi. Jism ko‘tarilishning eng baland nuqtasiga erishganda uning potensial energiyasi qanchaga ortadi?
\begin{enumerate}[label=\Alph*)]
\item 4,5
\item 18
\item 9
\item 12
\end{enumerate}

\item Har birining massasi 80 kg va qalinligi 0,2 m bo‘lgan uchta bir jinsli to‘g‘ri burchakli plita biri boshqasining yonida yer yuzasida gorizontal yotibdi. Plitalarni bir-birining ustiga terib qo‘yish uchun qanday minimal ish bajarish kerak? $g = 10 m/s^2$.
\begin{enumerate}[label=\Alph*)]
\item 240
\item 480
\item 960
\item 720
\end{enumerate}

\item Uzunligi 1,5 m va massasi 10 kg bo‘lgan yupqa lom gorizontal yuzada yotibdi. Uni yerga vertikal holatda qo‘yish uchun qanday minimal ish bajarish kerak? $g = 10 m/s^2$.
\begin{enumerate}[label=\Alph*)]
\item 150
\item 37,5
\item 75
\item 100
\end{enumerate}

\item Agar 200 kg massali brusokni gorizontal holatdan vertikal holatga o‘tkazilsa, uning potensial energiyasi qanchaga o‘zgaradi? Brusok 1 m uzunlikka va tomoni 20 cm bo‘lgan kvadrat kesimga ega. $g = 10 m/s^2$.
\begin{enumerate}[label=\Alph*)]
\item 400
\item 800
\item 1600
\item 1200
\end{enumerate}

\item Gorizontal tekislikda uzunligi 2 m va massasi 5 kg bo‘lgan yupqa zanjir yotibdi. Zanjirni bir uchidan ushlab, uning pastki uchi tekislikdan o‘zining uzunligiga teng bo‘lgan balandlikka ko‘tarishda bajariladigan minimal ish qanchaga teng? $g = 10 m/s^2$.
\begin{enumerate}[label=\Alph*)]
\item 75
\item 300
\item 200
\item 150
\end{enumerate}

\item Yuzasi 10 $m^2$ va chuqurligi 2 m bo‘lgan basseynni yarmigacha to‘ldirib turgan suvni yer yuzasiga chiqarib tashlash uchun qanday minimal ish (kJ da) bajarish kerak? $g = 10 m/s^2$.
\begin{enumerate}[label=\Alph*)]
\item 150
\item 75
\item 300
\item 200
\end{enumerate}

\item Yuzasi 15 $m^2$ bo‘lgan chuqur basseyn 1 m chuqurlikkacha suv bilan to‘ldirilgan va teng yarmidan vertikal to‘siq bilan to‘silgan. To‘siqni asta-sekinlik bilan shunday ko‘chiriladiki, bunda to‘siq basseynni 1:3 nisbatda bo‘ladi. Shunda qanday ish (kJ da) bajariladi? Suv to‘siq orasidan oqib o‘tmaydi. $g = 10 m/s^2$.
\begin{enumerate}[label=\Alph*)]
\item 12,5
\item 50
\item 25
\item 37,5
\end{enumerate}

\item 10 kN/m bikrlikli siqilgan prujina 50 J ga teng potensial energiya zahirasiga ega. Prujina necha santimetrga siqilgan?
\begin{enumerate}[label=\Alph*)]
\item 5
\item 10
\item 20
\item 15
\end{enumerate}

\item Agar 30 N kuch ta’siri ostida prujina 1 cm ga siqilsa, uni 20 cm ga siqish uchun bajarilishi kerak bo‘lgan ishni (J) toping.
\begin{enumerate}[label=\Alph*)]
\item 30
\item 120
\item 90
\item 60
\end{enumerate}

\end{enumerate}

\begin{table}[h!]
\centering
\caption{To‘g‘ri javoblar jadvali (1-27)}
\begin{tabular}{|c|c|c|c|c|c|c|c|c|c|}
\hline
№ & J. & № & J. & № & J. & № & J. & № & J. \\ \hline
1 & C & 7 & C & 13 & B & 19 & C & 25 & C \\ \hline
2 & A & 8 & B & 14 & C & 20 & B & 26 & B \\ \hline
3 & B & 9 & A & 15 & B & 21 & C & 27 & D \\ \hline
4 & D & 10 & D & 16 & A & 22 & B & & \\ \hline
5 & A & 11 & C & 17 & D & 23 & D & & \\ \hline
6 & C & 12 & A & 18 & B & 24 & A & & \\ \hline
\end{tabular}
\end{table}